\subsection{New Junta Theorems}
Next, we study connections between Talagrand-like quantities such as $\Expect{x}{\sqrt{s_f(x)}}$, and the notion of juntas.
For a point $x\in \{0,1\}^n$ and a subset $J\subseteq [n]$, we denote by $x_J$ the point in $\{0,1\}^J$ corresponding to
the value of $x$ on coordinates $J$.
\begin{definition}
  A function $f\colon\{0,1\}^n\to\{0,1\}$ is called a $j$-junta if there exists $J\subseteq [n]$ of size at most $j$,
  and $g\colon \{0,1\}^J\to\{0,1\}$ such that $f(x) = g(x_J)$ for all $x\in\{0,1\}^n$.

  We say a function $f\colon\{0,1\}^n\to\{0,1\}$ is $\eps$-close to a $j$-junta if there is a $j$-junta $g\colon\{0,1\}^n\to\{0,1\}$
  such that $\Prob{x}{f(x)\neq g(x)}\leq \eps$.
\end{definition}
On the one hand, it is possible that the quantity $\Expect{x}{\sqrt{s_f(x)}}$ is constant while the function $f$ is very far from being a junta,
as can be seen by taking $f$ to be the majority function. On the other hand, if we look at the expected sensitivity -- that is
$\Expect{x}{s_f(x)}$ -- a well known result of Friedgut~\cite{Friedgut} asserts that if $\Expect{x}{s_f(x)} = O(1)$, then $f$
is close to junta. More precisely, Friedgut's result asserts that if $\Expect{x}{s_f(x)} \leq A$, then
for all $\eps>0$, $f$ is $\eps$-close to a $2^{O(A/\eps)}$-junta. In light of this contrast, it is interesting to ask
what happens if an intermediate moment of the sensitivity, say the $p$th moment for $p\in (1/2,1)$, is constant.

We prove a strengthening of Friedgut's junta theorem~\cite{Friedgut} addressing this case, morally showing that for any $p>1/2$, bounded $p$-moment of
the sensitivity implies closeness to junta. More precisely:
\begin{thm}\label{thm:super_fried}
  For all $\eta>0$ there is $C = C(\eta)>0$ such that  the following holds.
  For all $\eps>0$ and $A>0$, and $1/2 + \eta\leq p\leq 1$, if $f\colon\{0,1\}^n\to\{0,1\}$ is a function such that
  $\Expect{x}{s_f(x)^{p}}\leq A$, then $f$ is $\eps$-close to a $J$-junta, where $J = 2^{C\left(A/\eps\right)^{1/p}}$.
\end{thm}
Theorem~\ref{thm:super_fried} is tight, and can be seen as a generalization of Friedgut's Junta Theorem (corresponding to the case that
$p=1$) that is somewhat close in spirit to Bourgain's Junta Theorem. Indeed, for all $p\geq 1/2$ the tribes function shows
this to be tight. Here, the tribes function is the function in which we take a partition $P_1\cup\ldots\cup P_m$ of $[n]$
and define ${\sf Tribes}(x) = 1$ if there is $i$ such that $x_{P_{i}} = \vec{1}$, and ${\sf Tribes}(x) = 0$ otherwise.
We take $P_i$ of equal size and the probability that ${\sf Tribes}(x)=1$ is roughly $1/2$; the correct choice
of parameters is $\card{P_i} = \log n - \log\log n + c$
and $m =  \frac{n}{\log n - \log\log n + c}$ where $c = O(1)$.
For this function one has $\Expect{x}{s_{{\sf Tribes}}(x)^{p}} = \Theta((\log n)^p)$. Indeed, as $s_{{\sf Tribes}}(x) = \Theta(\log n)$ with
constant probability we have the lower bound, and the upper bound follows by Jensen's inequality and the fact that the total influence of ${\sf Tribes}$
is $\Theta(\log n)$.
